\chapter{Học liên tục và monitoring}

\section{Chiến lược học liên tục}

Hệ thống thu dữ liệu mới qua feedback UI. Người dùng có thể sửa description, assignee, due date, đánh dấu false positive hoặc bổ sung missing task. Các correction được lưu trong PostgreSQL và export JSONL để phục vụ đánh giá hoặc huấn luyện. Trước khi training, dữ liệu được scrub PII, validate schema, loại trùng và chia train/validation/test theo meeting.

Retraining được thiết kế theo hướng có kiểm soát. Candidate model không tự động thay thế baseline sau mỗi feedback; nó được đánh giá bằng cùng benchmark và quality gates. Promotion manifest ghi lại model path, metrics và điều kiện pass/fail để giữ audit trail.

\section{Chỉ số monitoring}

Bảng \ref{tab:monitoring} định nghĩa các chỉ số vận hành cần theo dõi trong Prometheus/Grafana. Cách thiết kế monitoring và promotion gate dựa trên nguyên tắc MLOps có audit trail, đánh giá định kỳ và triển khai có kiểm soát \cite{shankar2022mlops}.

\begin{table}[H]
\centering
\footnotesize
\caption{Các metric monitoring}
\label{tab:monitoring}
\begin{tabularx}{\textwidth}{@{}>{\raggedright\arraybackslash}p{6.7cm}X@{}}
\toprule
\textbf{Chỉ số} & \textbf{Mục đích} \\
\midrule
\codefield{meeting_stage_duration_seconds} & Theo dõi độ trễ theo từng stage xử lý. \\
\codefield{meeting_llm_calls_total} và token metrics & Theo dõi số lần gọi LLM và kích thước prompt/response. \\
\codefield{meeting_jobs_total} & Theo dõi job hoàn tất và job lỗi. \\
\codefield{meeting_tasks_extracted_total} & Theo dõi số lượng task theo thời gian. \\
\codefield{meeting_hallucination_flags_total} & Theo dõi tín hiệu từ guardrails. \\
\codefield{meeting_anomaly_events_total} & Theo dõi anomaly detection bằng rolling Z-score. \\
\bottomrule
\end{tabularx}
\end{table}

\section{Suy giảm hiệu năng và drift}

Suy giảm hiệu năng được phát hiện bằng hai lớp. Online monitoring theo dõi latency, schema failures, guardrail flags, job failures và task count. Offline evaluation chạy benchmark định kỳ để so sánh precision, recall, F1, assignee accuracy và hallucination rate giữa baseline và candidate. Drift có thể đến từ domain meeting mới, cách nói mới, ngôn ngữ mới, chất lượng audio khác hoặc roster/alias thay đổi. Biện pháp giảm thiểu gồm versioning dataset/prompt/model, fixed benchmark, promotion gate và rollback bằng model manifest.
